\subsection{Summary of some works on RLWE/PLWE}\label{appendix:summary}

\subsubsection{WORK: Weak Instances of PLWE (https://eprint.iacr.org/2014/784.pdf)}

\textbf{Context and Problem:} \\
The paper analyzes the decisional-PLWE (Polynomial LWE) hardness assumption (which was introduced as a foundation for FHE schemes). The authors highlight a fundamental theoretical limitation: the worst-case to average-case reduction from the SVP to PLWE \emph{is only valid and proven for 2-power cyclotomic fields}. %When considering more general number rings, a distortion problem arises in the error distribution when transitioning from the continuous complex space (Minkowski embedding) to coefficient-wise sampling.

\textbf{Technique (guessing):} \\
The paper presents a targeted attack against PLWE in carefully constructed examples of number fields, rather than those generated by roots of unity. For a degree $n$ number field defined by $K = \mathbb{Q}[x]/(f(x))$ and a modulus $q$, if $f(1) \equiv 0 \pmod q$, the attack proves to be highly efficient. The technique is based on \emph{``guessing''} one of the $q$ possibilities for the value of the secret polynomial evaluated at $1$. This allows the attacker to distinguish PLWE samples (where error vectors are sampled coefficient-wise) with non-negligible probability in time $\tilde{\mathcal{O}}(q)$, vastly outperforming standard lattice attacks that require exponential time $\mathcal{O}(2^n)$.

\textbf{Main Contribution and Impact on Real Instances:} \\
Although the paper demonstrates that it is unsafe to rely directly on the PLWE assumption in arbitrary number fields, the authors explicitly clarify that \textbf{this does not constitute a practical attack and does not compromise real instances} of existing cryptosystems. This is primarily due to two reasons:
\begin{itemize}
    \item \textit{Choice of field:} Practical systems are based on 2-power cyclotomic rings, $R = \mathbb{Z}[x]/(\Phi_m(x))$. For these fields, $f(1) = \Phi_m(1) \not\equiv 0 \pmod q$, rendering the vulnerability completely inapplicable.
    \item \textit{Computational complexity:} The attack runs in $\mathcal{O}(q)$ time. In real-world FHE implementations, the modulus $q$ is cryptographically massive (e.g., $q \approx 2^{128}$), making an attack linear in $q$ entirely unfeasible in practice.
\end{itemize}



\subsubsection{WORK: Provably Weak Instances of Ring-LWE (https://eprint.iacr.org/2015/106.pdf)}

extension of the previous work

\textbf{Context and Problem:} \\
Building upon the foundational weaknesses in PLWE identified by Eisenträger, Hallgren, and Lauter (previous work), this paper addresses a major gap: EHL found weak number fields for PLWE, but failed to find fields that satisfied both the attack conditions and the security reductions to RLWE. This work investigates whether the mathematical structure of RLWE itself can be exploited, proving that the hardness of RLWE is highly sensitive to the specific properties of the chosen number field, the modulus $q$, and the polynomial basis.

\textbf{Methodology and Error Distortion:} \\
The core attack exploits scenarios where the defining polynomial $f(x)$ has a root of small order modulo $q$. The authors achieve this by mapping RLWE instances directly into PLWE instances. The success of this attack hinges on the \textit{spectral norm} of this transformation map, which dictates how much the error vectors are distorted during the mapping. If the distortion is sufficiently small (governed by the ratio $\sqrt{q}/n$) the noise does not mask the secret effectively. The resulting algorithm runs in linear time with respect to the modulus, $\mathcal{O}(q)$.

\textbf{Main Contribution and Impact on Real Instances:} \\
This paper presents the first successful attacks on the Decision RLWE problem (and consequently the Search RLWE problem for Galois fields) for explicitly constructed families of number fields. The practical implications are significant:
\begin{itemize}
    \item \textit{Vulnerability of standard fields:} The authors prove that even mathematically secure 2-power cyclotomic fields can be completely broken if they are instantiated using an alternative (bad) polynomial basis.
    \item \textit{Practical implementation:} The attack is not just theoretical; it was successfully implemented at cryptographic sizes (e.g., $n=1024$ and $q = 2^{31}-1$ was broken in about 13 hours).
    \item \textit{Impact on real protocols:} While FHE schemes remain safe due to their massive moduli (128 to 512 bits, making an $\mathcal{O}(q)$ attack impossible), the paper warns that RLWE used in Key Exchange protocols (such as BCNS for TLS) often use smaller moduli (e.g., $q \approx 2^{32}$). These protocols are highly vulnerable if weak parameters are inadvertently chosen.
\end{itemize}


\subsubsection{WORK: Attacks on the Search RLWE Problem with Small Errors (https://epubs.siam.org/doi/epdf/10.1137/16M1096566)}

\textbf{Context and Problem:} \\
While RLWE has rigorous security proofs for specific parameter sets, practical implementations often use narrower error distributions to improve computational efficiency. Previous attacks exploited simple homomorphisms to finite fields $\mathbb{F}_q$, but they were limited to non-Galois fields, meaning they could only solve the Decision variant of RLWE and could not recover the actual secret key (the Search variant).

\textbf{Methodology and ``Subfield Vulnerabilities'':} \\
This paper pioneers a novel mathematical attack vector. Instead of mapping to $\mathbb{F}_q$, the authors use homomorphisms to higher-degree finite fields, $\mathbb{F}_{q^f}$ (where $f > 1$). To detect the hidden structure, they deploy the \textbf{statistical chi-square test} to distinguish the mapped error distribution from a uniform distribution. This technique exploits geometric ``subfield vulnerabilities'' (specifically how certain prime ideals behave as sublattices). For a number field of degree $n$ and residue degree $f$, the attack runs in time $\mathcal{O}(nq^{2f})$. 

\textbf{Main Contribution and Impact on Real Instances:} \\
The paper makes several major breakthroughs in RLWE cryptanalysis, though it explicitly bounds its impact on current real-world implementations:
\begin{itemize}
    \item \textit{Search-to-Decision Reduction:} The authors successfully prove a search-to-decision reduction for general Galois fields with unramified moduli, allowing their attack to fully recover secret keys.
    \item \textit{Attacks on Prime Cyclotomics:} They demonstrate devastatingly fast attacks on prime cyclotomic rings. For instance, in a cyclotomic ring $\mathbb{Q}(\zeta_{809})$ with dimension $n=808$ and modulus $q=809$ (the ramified prime), the decision RLWE problem is broken in just 35 seconds. 
    \item \textit{Impact on Practical Security:} The vulnerabilities apply to error widths that are narrower than the theoretically secure bounds (e.g., $r = \omega(\sqrt{\log n})$). Importantly, most practical HE schemes strictly use \textbf{2-power cyclotomic rings}, which are completely immune to this specific subfield attack. Therefore, existing secure implementations remain safe, but the paper serves as a severe warning against using prime cyclotomic fields or aggressively shrinking error parameters for efficiency.
\end{itemize}


\subsubsection{WORK: On the Ring-LWE and Polynomial-LWE Problems (https://eprint.iacr.org/2018/170.pdf)}

\textbf{Context and Problem:} \\
Prior to this work, the cryptographic landscape was fragmented into three distinct but related hardness assumptions: 
\begin{enumerate}
    \item Dual-RLWE (operating in the fractional dual ideal $\mathcal{O}_K^\vee$),
    \item Primal-RLWE (operating directly in the ring of integers $\mathcal{O}_K$),
    \item PLWE (operating on simple polynomial coefficients in $\mathbb{Z}[x]/\langle f\rangle$).
\end{enumerate}
While these problems trivially collapse into a single equivalent problem for 2-power cyclotomic fields, their relationship in arbitrary number fields remained an open and critical question, especially given the rising concern that specific fields might be weak against algebraic attacks.

\textbf{Methodology and Algebraic Tools:} \\
The authors prove that these three problems (in both their Search and Decision variants) are essentially equivalent for massive classes of rings by constructing polynomial-time reductions between them. The mathematical heavy lifting involves:
\begin{itemize}
    \item \textbf{The Conductor Ideal:} To bridge the gap between the formal ring of integers $\mathcal{O}_K$ and the easily computable polynomial order $\mathbb{Z}[\alpha]$, the authors introduce the \textit{conductor ideal} $\mathcal{C}_{\mathbb{Z}[\alpha]}$. This is a novel application in RLWE cryptanalysis.
    \item \textbf{Condition Number Bounds:} To prove that transitioning from PLWE (coefficient embedding) to RLWE (Minkowski embedding) does not stretch the error vectors exponentially, they rigorously bound the \textit{condition number} of the transformation map for large families of polynomials (e.g., $x^n + x \cdot P(x) - a$).
    \item \textbf{Oracle Hidden Center Problem (OHCP):} Leveraging techniques from Peikert et al., they construct a Search-to-Decision reduction for Dual-RLWE across arbitrary fields by mapping Gaussian combinations of RLWE samples to an OHCP instance.
\end{itemize}

\textbf{Main Contribution and Impact on Real Instances:} \\
The central thesis of the paper is a unifying one: \textbf{the six variants of RLWE/PLWE are qualitatively equivalent}. 
\begin{itemize}
    \item \textit{Theoretical Reassurance:} This equivalence strongly suggests that there is no fundamental design flaw in one variant that does not exist in the others. If PLWE is weak for a specific polynomial, Primal and Dual RLWE will also be weak for that generated field, and vice versa.
    \item \textit{Error Growth Constraints:} The reductions inherently cause the error rate to grow. This growth is bounded polynomially by factors like the degree $n$, the root discriminant $|\Delta_K|^{1/n}$, and the norm of the conductor ideal.
    \item \textit{Practical Cryptographic Design:} While Dual-RLWE is mathematically purer, it requires computing the ring of integers $\mathcal{O}_K$ (which is computationally hard for arbitrary fields). This paper justifies the engineering preference for PLWE, as it avoids computing $\mathcal{O}_K$, though decoding the noise in PLWE still requires finding a good lattice basis for $\mathbb{Z}[\alpha]$.
\end{itemize}


\subsubsection{EXTRA: Recovering Short Generators of Principal Ideals in Cyclotomic Rings (https://eprint.iacr.org/2015/313.pdf)}




\subsubsection{EXTRA: On the Security of the Multivariate Ring Learning with Errors Problem (https://eprint.iacr.org/2018/966.pdf)}

\textbf{Context and Problem:} \\
In an attempt to make lattice-based cryptography even more efficient (specifically for multidimensional signal and image processing), Pedrouzo-Ulloa et al. introduced the \textbf{Multivariate Ring Learning with Errors (m-RLWE)} problem. Instead of working in a univariate polynomial ring $\mathbb{Z}_q[x]/(f(x))$, m-RLWE operates over multivariate rings like $\mathbb{Z}_q[x, y]/(f(x), g(y))$. The original authors proposed using tensor products of power-of-two cyclotomic rings, incorrectly claiming that the security of the system scaled with the \textit{product} of the individual degrees (e.g., $n_1 \times n_2$).

\textbf{Methodology and Attack Vector:} \\
This paper systematically dismantles the security claims of m-RLWE. The authors identify a fatal flaw in the assumption that the tensor product of cyclotomic rings creates a new, mathematically independent, high-dimensional field. Because the defining cyclotomic polynomials share algebraically related roots, the multivariate ring is actually highly redundant. 
The attack leverages ``smallness-preserving'' ring homomorphisms to project the high-dimensional m-RLWE problem down into multiple standard univariate RLWE problems. Crucially, the dimension of these sub-problems is only the \textit{maximum} of the individual degrees (e.g., $\max\{n_1, n_2\}$), not their product. Furthermore, the noise during this projection only grows by a tiny factor (the square root of the other degree, $\sqrt{n_2}$).

\textbf{Main Contribution and Impact on Real Instances:} \\
The paper's findings are devastating to the original m-RLWE proposal and serve as a severe warning against naive algebraic optimizations:
\begin{itemize}
    \item \textit{Shattered Security Estimates:} Parameters that the original authors estimated to offer over 2600 bits of security (and up to 146,000 bits) were proven to actually offer only 32 to 98 bits of security.
    \item \textit{Practical Implementation:} The authors successfully implemented their attack. A proposed parameter set equivalent to solving a Bounded Distance Decoding (BDD) problem in a massive lattice of dimension 16,384 was completely broken in less than two weeks on standard hardware (and just a few hours on 128 cores).
    \item \textit{Warning for Module-LWE:} The authors generalize their warning to Module-LWE (the foundation of standards like NIST's ML-KEM/Kyber). They demonstrate that attempting to optimize Module-LWE by naively introducing extra ring structures (treating the module inner product as a polynomial product) often results in a total breakdown of security.
\end{itemize}